\documentclass[aspectratio=169]{beamer}
\usetheme{metropolis}
\usepackage{appendixnumberbeamer}
\usepackage{booktabs, hyperref}

\input{mgmt675-style}

\usepackage{tikz}
\usetikzlibrary{shapes.geometric, arrows.meta, positioning, calc}

% Title info
\subtitle{MGMT 675: Generative AI for Finance}
\title{Module 4: Claude Desktop and Claude Code}
\author{Kerry Back}
\date{}

\begin{document}

\maketitle

% ========================================
% SECTION 1: CLAUDE DESKTOP OVERVIEW
% ========================================
\section{Claude Desktop Overview}

\begin{frame}{Three Tabs: Chat, Cowork, Code}
\begin{columns}[T]
\begin{column}{0.32\textwidth}
\begin{shadedbox}[title=\textbf{Chat (Analysis)}]
\begin{itemize}\footnotesize
  \item Sandboxed Python in the browser
  \item No local file access
  \item Must upload data manually
\end{itemize}
\end{shadedbox}
\end{column}
\begin{column}{0.32\textwidth}
\begin{shadedbox}[title=\textbf{Cowork}]
\begin{itemize}\footnotesize
  \item Runs in a local VM
  \item Files synced to/from VM
  \item \alert{No internet access}
\end{itemize}
\end{shadedbox}
\end{column}
\begin{column}{0.32\textwidth}
\begin{shadedbox}[title=\textbf{Code}]
\begin{itemize}\footnotesize
  \item Runs on \alert{your machine}
  \item Full local file access
  \item Full internet access
\end{itemize}
\end{shadedbox}
\end{column}
\end{columns}
\vspace{0.3cm}
\begin{center}
\alert{Code mode} gives Claude direct access to your machine --- no sandbox, no VM, no upload step.
\end{center}
\end{frame}

\begin{frame}{Which Tool When?}
\begin{itemize}
  \item \textbf{Chat} for questions, explanations, and brainstorming
  \item \textbf{Artifacts} for quick interactive tools (e.g., a DCF calculator with sliders)
  \item Code can do everything Cowork can do, plus it has \alert{full internet access} and direct access to your local files
  \item Code can also build and deploy interactive tools (like Artifacts)
  \item Cowork's advantage is that it requires \textbf{no local software setup} and runs in an isolated sandbox
\end{itemize}
\end{frame}

\begin{frame}{The Class Installer}
The class installer sets up everything you need in one step:
\vspace{0.3cm}
\begin{itemize}
  \item \textbf{Python} --- language Claude Code uses to execute analysis
  \item \textbf{Claude Desktop} --- the application with built-in Claude Code
  \item \textbf{Skills} --- pre-packaged instructions for Excel, Word, PowerPoint, and more
\end{itemize}
\vspace{0.3cm}
Download the installer for your platform:
\begin{itemize}
  \item \href{https://mgmt675.kerryback.com/install/windows/windows-install.zip}{Windows Installer}
  \item \href{https://mgmt675.kerryback.com/install/mac/install.sh}{Mac Installer}
\end{itemize}
\end{frame}

\begin{frame}{Running the Installer}
\begin{columns}[T]
\begin{column}{0.48\textwidth}
\begin{shadedbox}[title=\textbf{Windows}]
\begin{enumerate}\small
  \item Download \texttt{windows-install.zip}
  \item Extract the contents
  \item Right-click \texttt{INSTALL.bat}\\$\rightarrow$ \textbf{Run as administrator}
\end{enumerate}
\end{shadedbox}
\end{column}
\begin{column}{0.48\textwidth}
\begin{shadedbox}[title=\textbf{Mac}]
\begin{enumerate}\small
  \item Download \texttt{install.sh} to Downloads
  \item Open \textbf{Terminal}
  \item Paste and run:\\[2pt]
    \texttt{\footnotesize cd \textasciitilde/Downloads \&\& bash install.sh}
\end{enumerate}
\end{shadedbox}
\end{column}
\end{columns}
\end{frame}

\begin{frame}{Getting Started}
\begin{enumerate}
  \item Open Claude Desktop $\rightarrow$ \textbf{Code} tab
  \item Select a working folder for your project
  \item Choose a model: \textbf{Sonnet} (fast) or \textbf{Opus} (thorough)
  \item Set permission mode to \textbf{auto-accept} (recommended for exercises)
\end{enumerate}
\end{frame}

\begin{frame}{CLAUDE.md}
Claude reads \texttt{CLAUDE.md} files automatically at the start of every session --- no need to repeat instructions.
\vspace{0.3cm}
\begin{itemize}
  \item \textbf{Global} (\texttt{\textasciitilde/.claude/CLAUDE.md}): applies to every project --- output preferences, preferred libraries, personal conventions
  \item \textbf{Per-Project} (project root): applies only to this folder --- data context, domain guidance, project-specific workflows
\end{itemize}
\vspace{0.3cm}
The installer pre-configured \textbf{skills} for Excel, Word, and PowerPoint.  Prompt Claude Code to document them in your global \texttt{CLAUDE.md}:
\vspace{0.2cm}

\textit{``Read the skills in my .claude/skills folder and add a summary of each to my global CLAUDE.md.''}
\vspace{0.2cm}

Once documented, Claude will \alert{automatically} follow those instructions without needing slash commands.
\end{frame}

\begin{frame}{Example}
\textit{``Create an Excel workbook with a loan amortization table for a \$300{,}000 mortgage at 6.5\% annual interest with a 30-year term and monthly payments.  Use formulas so I can change the loan amount, rate, and term.  Include columns for payment number, payment, principal, interest, and remaining balance.''}
\end{frame}

\section{Claude Code in the Terminal and VS Code}

\begin{frame}{Terminal Alternative}
The Code tab in Claude Desktop and the terminal command \texttt{claude} use the \alert{same engine}.
\vspace{0.3cm}
\begin{columns}[T]
\begin{column}{0.5\textwidth}
\begin{shadedbox}[title=\textbf{Claude Desktop (Code tab)}]
\begin{itemize}\small
  \item Visual interface with file browser
  \item Click to select working folder
  \item Good for getting started
\end{itemize}
\end{shadedbox}
\end{column}
\begin{column}{0.5\textwidth}
\begin{shadedbox}[title=\textbf{Terminal (\texttt{claude})}]
\begin{itemize}\small
  \item Type \texttt{claude} in any terminal
  \item Same skills, same config
  \item Preferred by power users
\end{itemize}
\end{shadedbox}
\end{column}
\end{columns}
\vspace{0.3cm}
Both approaches share the same skills, MCP servers, and configuration files.
\end{frame}

\begin{frame}{Installing the Terminal and VS Code}
\begin{columns}[T]
\begin{column}{0.5\textwidth}
\begin{shadedbox}[title=\textbf{Claude Code (Terminal)}]
\begin{itemize}\small
  \item \textbf{Mac/Linux:} Open Terminal and run:\\[2pt]
    \texttt{\footnotesize curl -fsSL https://claude.ai/install.sh | bash}
  \item \textbf{Windows:} Open PowerShell and run:\\[2pt]
    \texttt{\footnotesize irm https://claude.ai/install.ps1 | iex}
  \item Windows also requires \href{https://git-scm.com/downloads/win}{Git for Windows}
  \item Verify with \texttt{claude --version}
\end{itemize}
\end{shadedbox}
\end{column}
\begin{column}{0.5\textwidth}
\begin{shadedbox}[title=\textbf{VS Code + Extensions}]
\begin{enumerate}\small
  \item Download \href{https://code.visualstudio.com}{VS Code} and install
  \item Open Extensions sidebar (\texttt{Ctrl+Shift+X} on Windows, \texttt{Cmd+Shift+X} on Mac)
  \item Search \textbf{Python} (by Microsoft) $\rightarrow$ Install
  \item Search \textbf{Claude Code} (by Anthropic) $\rightarrow$ Install
\end{enumerate}
\end{shadedbox}
\end{column}
\end{columns}
\vspace{0.3cm}
Claude Code updates automatically.  The VS Code extension uses the same engine as the terminal --- install either or both.
\end{frame}

\begin{frame}{Using Claude Code in the Terminal and VS Code}
\begin{columns}[T]
\begin{column}{0.5\textwidth}
\begin{shadedbox}[title=\textbf{Terminal}]
\begin{itemize}\small
  \item \texttt{cd} into your project folder, then type \texttt{claude}
  \item First launch: sign in with your Anthropic account
  \item Chat naturally --- Claude reads and writes files in that folder
  \item Type \texttt{/help} for commands, \texttt{/model} to switch models
  \item Press \texttt{Esc} to interrupt, \texttt{Ctrl+C} twice to exit
\end{itemize}
\end{shadedbox}
\end{column}
\begin{column}{0.5\textwidth}
\begin{shadedbox}[title=\textbf{VS Code}]
\begin{itemize}\small
  \item Open your project folder in VS Code
  \item Click the Claude Code icon in the sidebar (or \texttt{Ctrl+Shift+P} $\rightarrow$ ``Claude Code: Open'')
  \item Sign in on first use
  \item Claude can see your open files, file explorer, and terminal
  \item Code changes appear as inline diffs you can accept or reject
\end{itemize}
\end{shadedbox}
\end{column}
\end{columns}
\vspace{0.3cm}
\alert{Both share the same skills and configuration.}  VS Code adds a file explorer and visual diffs; the terminal is lightweight and fast.
\end{frame}

% ========================================
% SECTION 3: WORKING WITH MULTIPLE FILES
% ========================================
\section{Working with Multiple Files}

\begin{frame}{Merging Data from Different Sources}
Real financial data rarely lives in a single file or format.  Claude Code can read CSV, Excel, PDF, and Word files, then \alert{combine them into a unified dataset}.
\vspace{0.3cm}
\begin{center}
\begin{tikzpicture}[
  node distance=0.6cm and 0.8cm,
  every node/.style={font=\small},
  sbox/.style={draw=titlegray, rounded corners, minimum width=1.8cm, minimum height=0.7cm, fill=excelinput, text=titlegray, font=\scriptsize\bfseries, align=center},
  sarrow/.style={-{Stealth[length=2.5mm]}, thick, draw=accentblue}
]
\node[sbox] (csv) {CSV};
\node[sbox, below=0.4cm of csv] (xlsx) {Excel};
\node[sbox, below=0.4cm of xlsx] (pdf) {PDF};

\node[sbox, right=1.5cm of xlsx, fill=accentblue!15, minimum width=2.2cm] (claude) {Claude\\Code};
\node[sbox, right=1.5cm of claude, fill=alertorange!15, minimum width=2.2cm] (unified) {Unified\\Table};

\draw[sarrow] (csv.east) -- (claude.west |- csv.east);
\draw[sarrow] (xlsx) -- (claude);
\draw[sarrow] (pdf.east) -- (claude.west |- pdf.east);
\draw[sarrow] (claude) -- (unified);
\end{tikzpicture}
\end{center}
\vspace{0.3cm}
\textbf{Common scenarios:}
\begin{itemize}\small
  \item Holdings from multiple brokerages (xlsx + csv)
  \item Loan tape (CSV) + appraisal report (PDF) + policy memo (DOCX)
  \item Quarterly reports from divisions with different column names
\end{itemize}
\end{frame}

\begin{frame}{What Claude Handles}
\begin{itemize}
  \item Reads each file format natively
  \item Reconciles column names and data types
  \item Flags mismatches and missing values
  \item Outputs a clean, unified table
\end{itemize}
\vspace{0.5cm}
\alert{You describe what should be merged and how.  Claude handles the parsing and code.}
\end{frame}

\begin{frame}{Filtering and Aggregating}
Once data is combined, the next step is usually filtering, sorting, or grouping.
\vspace{0.3cm}
\begin{columns}[T]
\begin{column}{0.5\textwidth}
\begin{shadedbox}[title=\textbf{Filtering}]
\begin{itemize}\small
  \item ``Keep only loans above \$1M''
  \item ``Filter to stocks with P/E > 0''
  \item ``Exclude rows with missing revenue''
\end{itemize}
\end{shadedbox}
\end{column}
\begin{column}{0.5\textwidth}
\begin{shadedbox}[title=\textbf{Aggregating}]
\begin{itemize}\small
  \item ``Average revenue by sector and quarter''
  \item ``Sum holdings by asset class''
  \item ``Rank funds by Sharpe ratio''
\end{itemize}
\end{shadedbox}
\end{column}
\end{columns}
\vspace{0.3cm}
State your intent in plain English.  Claude writes the pandas code, runs it, and shows you the result.  Always \alert{check the filter logic} --- this is where silent errors happen most often (Module~5).
\end{frame}

\begin{frame}{Example}
Download \href{https://kerryback.com/mgmt675/files/aggregation.zip}{aggregation.zip} and extract the workbooks into a folder.
\vspace{0.3cm}
Each Excel file contains a data table.  Column names vary across files and some columns are missing.  Ask Claude Code to:
\begin{enumerate}
  \item Summarize the columns present, row count, and basic statistics for each file
  \item Combine all files into a single table, reconciling the varying column names
  \item Save the result as \texttt{combined.xlsx}
\end{enumerate}
\vspace{0.3cm}
This is the same ``divide and conquer'' pattern used in professional data pipelines --- read each source independently, then unify.
\end{frame}

\end{document}
